\documentclass[11pt]{article}
\usepackage[margin=1in]{geometry}
\usepackage{amsmath,amssymb,amsfonts,mathtools}
\usepackage{array,booktabs,enumitem,graphicx,xcolor,hyperref}
\usepackage[T1]{fontenc}
\usepackage[utf8]{inputenc}
\title{World Sheet and Space Time Physics in Two Dimensional (Super) String
  Theory}
\author{Source-backed arXiv sample}
\date{}
\begin{document}
\maketitle
\begin{abstract}
We show that tree level ``resonant'' \$N\$ tachyon scattering amplitudes, which define a sensible ``bulk'' S -- matrix in critical (super) string theory in any dimension, have a simple structure in two dimensional space time, due to partial decoupling of a certain infinite set of discrete states. We also argue that the general (non resonant) amplitudes are determined by the resonant ones, and calculate them explicitly, finding an interesting analytic structure. Finally, we discuss the space time interpretation of our results.
\end{abstract}
\section{Source Metadata}
This sample is based on arXiv record \texttt{hep-th/9109005}. Authors: P. Di Francesco and D. Kutasov.
\section{Extracted Prose 1}
We show that tree level ``resonant'' math expression tachyon scattering amplitudes, which define a sensible ``bulk'' S -- matrix in critical (super) string theory in any dimension, have a simple structure in two dimensional space time, due to partial decoupling of a certain infinite set of discrete states. We also argue that the general (non resonant) amplitudes are determined by the resonant ones, and calculate them explicitly, finding an interesting analytic structure. Finally, we discuss the space time interpretation of our results.
\section{Extracted Prose 2}
E. Brezin, V. Kazakov, and Al. Zamolodchikov, Nucl. Phys. B338 (1990) 673; D. Gross and N. Miljkovi\textbackslash{}'c, Phys. Lett 238B (1990) 217; P. Ginsparg and J. Zinn-Justin, Phys. Lett 240B (1990) 333; G. Parisi, Phys. Lett. 238B (1990) 209. , it was understood that the center of mass in these 2D string theories is described by free fermion quantum mechanics T. Banks, M. Douglas, N. Seiberg and S. Shenker, Phys. Lett. 238B (1990) 279. S. Das and A. Jevicki, Mod. Phys. Lett. A5 (1990) 1639.
\end{document}
